% ===== Opening: title, hook, the book and the series, today's plan (~5 min) =====

\webinartitle{8}{Գլուխ 8. Երբ ուղեղը «փչանում» է}%
  {\chip{raincloud}{դեպրեսիա} \chip{battery}{սթրես} \chip{monkey}{կախվածություն} \chip{speakinghead}{հալյուցինացիաներ} \chip{party}{վերջաբան}}%
  {confused}{Տեղի է ունեցել անսպասելի սխալ}{{raincloud}{sunbehindcloud}{rainbow}}

% Hook: the chapter's opening list of everything the brain gets wrong, recapping the series
\begin{frame}{Ի՞նչ ենք իմացել ուղեղի մասին մինչ այժմ}
  \begin{columns}[c]
    \begin{column}{0.56\textwidth}
      \begin{popup}{սխալներ.log}
        \small
        \begin{enumerate}\setlength{\itemsep}{0.2mm}
          \item Խմբագրում է հիշողություններդ
          \item Վախենում է անվնաս բաներից
          \item Խառնում է սնունդդ, քունդ, շարժումդ
          \item Համոզում է, որ փայլուն ես
          \item Ընկալածիդ կեսը «հորինում» է
          \item Զգացմունքներից անտրամաբանական է դառնում
          \item Արագ ընկերանում է, արագ թշնամանում
        \end{enumerate}
      \end{popup}
    \end{column}
    \begin{column}{0.4\textwidth}
      \onslide<2->{%
      \begin{chatbox}
        Ո՞ր համարն է քեզ հետ ամենահաճախը։ Գրիր միայն թիվը։
      \end{chatbox}}
      \vspace{1mm}
      \onslide<3->{%
      \begin{card}[colback=brainpink!30!cardwhite]
        Ու այս ամենը՝ երբ ուղեղը \textbf{նորմալ} է աշխատում։ Իսկ եթե այն իսկապես «փչանա՞»
      \end{card}}
    \end{column}
  \end{columns}
\end{frame}

% The book and the series, for newcomers
\begin{frame}{Գիրքը և շարքը}
  \begin{columns}[T]
    \begin{column}{0.41\textwidth}
      \begin{card}
        {\blackfont\Large The Idiot Brain} {\small (2016)}\par\smallskip
        \textbf{Դին Բըրնեթ}՝ նեյրոգիտնական։ Դասավանդել է Քարդիֆի համալսարանի հոգեբուժության ու կլինիկական նեյրոգիտության ինստիտուտում։\par\smallskip
        Այսօրվա թեման նրա մասնագիտությունն է։
      \end{card}
    \end{column}
    \begin{column}{0.55\textwidth}
      {\blackfont\small 8 գլուխ, 8 վեբինար}\par\smallskip
      {\color{plum!60!cream}\footnotesize
      \begin{tabular}{@{}l@{\hspace{4mm}}l@{}}
        1. Մարմնի կառավարում & 5. Ընկալում\\
        2. Հիշողություն & 6. Անհատականություն\\
        3. Վախ & 7. Ուրիշ մարդիկ\\
        4. Խելք &
      \end{tabular}}\par\medskip
      \chip[brainpink]{brain}{Այսօր՝ 8. Երբ ուղեղը «փչանում» է}\par\smallskip
      \chip[butter]{book}{+ վերջաբանը}\par\medskip
      {\small Շարքի \textbf{վերջին} վեբինարն է. վերջում ամբողջ գիրքը կնայենք մեկ էջում։}
    \end{column}
  \end{columns}
\end{frame}

% Plan for today
\begin{frame}{Այսօր}
  \begin{columns}[T]
    \begin{column}{0.27\textwidth}
      \begin{card}[colback=butter!45!cream]
        {\blackfont Ներածություն}\par\smallskip
        Ի՞նչն է «նորմալ» ուղեղը, ու ո՞վ է դա որոշում
      \end{card}
    \end{column}
    \begin{column}{0.31\textwidth}
      \begin{card}[colback=lavender!45!cream]
        {\blackfont Մինչ ընդմիջումը}\par\smallskip
        \chip{raincloud}{դեպրեսիա}\par\smallskip
        \chip{battery}{սթրես ու պոռթկում}
      \end{card}
    \end{column}
    \begin{column}{0.36\textwidth}
      \begin{card}[colback=mint!45!cream]
        {\blackfont Ընդմիջումից հետո}\par\smallskip
        \chip{monkey}{կախվածություն}\par\smallskip
        \chip{speakinghead}{հալյուցինացիաներ}\par\smallskip
        \chip{party}{վերջաբան}
      \end{card}
    \end{column}
  \end{columns}
  \vspace{3mm}
  \centering
  {\small \emoji{coffee}\ Մեջտեղում՝ 5 րոպե ընդմիջում \qquad \emoji{speech}\ Հարցերը գրիր չաթում ցանկացած պահի}
\end{frame}
